
\documentclass[11pt]{article}
\usepackage{amsmath}
\begin{document}
\section*{BernoulliSolver Command-Line Test Cases}

This report summarizes five command-line validation simulations executed with \texttt{Bending\_Torsion\_Beam\_Test\_Case/BernoulliSolver}.  Each simulation was launched using an mdpa mesh, a material json file, and command-line boundary-condition/load options.  No testcase-specific ProjectParameters file was used for the command-line runs.

\subsection*{Material and Section Data}
The tests use
\[
E = 206.9\times 10^9,\qquad \nu=0.29,\qquad G=\frac{E}{2(1+\nu)},\qquad A=I_z=J=1.0,\qquad L=10.0.
\]
The load convention is global: \texttt{--load node fx fy fz mx my mz}.

\subsection*{Analytical References}
\begin{itemize}
\item Cantilever point load: $v(L)=PL^3/(3EI_z)$.
\item Simply supported center point load: $v(L/2)=PL^3/(48EI_z)$.
\item Simply supported UDL: $v(L/2)=5qL^4/(384EI_z)$.
\item Pure torsion: $\theta_x(L)=TL/(GJ)$.
\item Combined loading: superposition of $\theta_x(L)=TL/(GJ)$, $v_y(L)=M_zL^2/(2EI_z)$, and $\theta_z(L)=M_zL/(EI_z)$.
\end{itemize}

\subsection*{Numerical Comparison}
\begin{table}[h]
\centering
\scriptsize
\begin{tabular}{llrrr}
\hline
Case & Quantity & Numerical & Analytical & Error (\%) \\
\hline
Cantilever end point load & tip Uy & -1.611084300000e-06 & -1.611084259707e-06 & 2.501e-06 \\
Simply supported center point load & center Uy & -1.006927700000e-07 & -1.006927662317e-07 & 3.742e-06 \\
Simply supported uniform distributed load & center Uy & -6.293298000000e-07 & -6.293297889480e-07 & 1.756e-06 \\
Fixed-free torsion & tip Rx & 1.246979200000e-07 & 1.246979217013e-07 & 1.364e-06 \\
Combined bending and torsion & tip Rx & 1.246979200000e-07 & 1.246979217013e-07 & 1.364e-06 \\
Combined bending and torsion & tip Uy & 1.208313200000e-07 & 1.208313194780e-07 & 4.320e-07 \\
Combined bending and torsion & tip Rz & 2.416626400000e-08 & 2.416626389560e-08 & 4.320e-07 \\
\hline
\end{tabular}
\end{table}

\subsection*{Qualitative Evaluation}
The cantilever point-load case shows the expected downward cubic bending shape with zero fixed-end displacement and rotation.  The simply supported point-load and UDL cases show zero transverse displacement at the supports and maximum downward displacement at midspan.  The torsion case has negligible translational displacement and a nonzero twist about the beam axis.  The combined loading case reproduces linear superposition of torsion and pure end bending.

\subsection*{Output Structure}
The five VTK output folders are:
\begin{itemize}
\item \texttt{Bending\_Torsion\_Beam\_Test\_Case/BernoulliSolver\_TestCases/01\_cantilever\_point\_load/BeamModelPart\_0\_0.vtk}
\item \texttt{Bending\_Torsion\_Beam\_Test\_Case/BernoulliSolver\_TestCases/02\_simply\_supported\_point\_load/BeamModelPart\_0\_0.vtk}
\item \texttt{Bending\_Torsion\_Beam\_Test\_Case/BernoulliSolver\_TestCases/03\_simply\_supported\_udl/BeamModelPart\_0\_0.vtk}
\item \texttt{Bending\_Torsion\_Beam\_Test\_Case/BernoulliSolver\_TestCases/04\_torsion/BeamModelPart\_0\_0.vtk}
\item \texttt{Bending\_Torsion\_Beam\_Test\_Case/BernoulliSolver\_TestCases/05\_combined\_bending\_torsion/BeamModelPart\_0\_0.vtk}

\end{itemize}

\subsection*{Conclusion}
All five command-line simulations agree with the corresponding Euler--Bernoulli analytical solutions to VTK ASCII output precision.  The maximum reported percentage errors are on the order of $10^{-6}$ percent for the scalar quantities checked, indicating that the stiffness assembly, local-to-global transformation, command-line boundary conditions, command-line loads, and VTK displacement/rotation output are consistent for these benchmark cases.
\end{document}
